\documentclass[tikz,border=12pt]{standalone}
\usepackage[T1]{fontenc}
\usepackage{mathpazo}
\usepackage{amsmath,amssymb}
\usetikzlibrary{arrows.meta, positioning, calc}

\definecolor{stblue}{HTML}{7B97AD}
\definecolor{rwgold}{HTML}{B8995E}
\definecolor{sage}{HTML}{7A9A7E}
\definecolor{dkbrown}{HTML}{3D3530}
\definecolor{wmgray}{HTML}{7B7067}
\definecolor{cream}{HTML}{F9F6F0}
\definecolor{border}{HTML}{D4C9B8}
\definecolor{terracotta}{HTML}{C47A6A}
\definecolor{lavender}{HTML}{907EA0}

\begin{document}
\begin{tikzpicture}[
  every node/.style={text=dkbrown, font=\large},
  >=Stealth[length=4mm, width=3mm],
]

% Background
\fill[cream, rounded corners=14pt] (-1.5,-3) rectangle (16,11);
\draw[border, line width=1.5pt, rounded corners=14pt] (-1.5,-3) rectangle (16,11);

% Title
\node[font=\LARGE\bfseries, text=dkbrown] at (7.25,10.2) {Nonconvex PL Function and Gradient Domination};

% Plot area: x in [-5,5] -> tikz [0,14], y in [0,50] -> tikz [0,8.5]
% x_tikz = (x+5)/10 * 14 = (x+5)*1.4
% y_tikz = y/50 * 8.5 = y*0.17

% Axes
\draw[->, line width=1pt, dkbrown] (0,0) -- (14.5,0) node[right, font=\large] {$x$};
\draw[->, line width=1pt, dkbrown] (7,0) -- (7,9.2) node[above, font=\large] {$y$};

% x-axis tick labels
\foreach \x/\lab in {0/{$-5$}, 2.8/{$-3$}, 5.6/{$-1$}, 7/0, 8.4/1, 11.2/3, 14/5} {
  \node[font=\large, text=wmgray, below] at (\x,-0.3) {\lab};
}

% y-axis tick labels
\foreach \v/\ytik in {10/1.7, 20/3.4, 30/5.1, 40/6.8} {
  \node[font=\large, text=wmgray, left] at (6.8,\ytik) {\v};
  \draw[wmgray, line width=0.5pt] (6.85,\ytik) -- (7.15,\ytik);
}

% Grid lines
\foreach \ytik in {1.7, 3.4, 5.1, 6.8} {
  \draw[wmgray, opacity=0.12, line width=0.5pt] (0,\ytik) -- (14,\ytik);
}

% f(x) = x^2 + 3*sin^2(x)
% f'(x) = 2x + 3*sin(2x)
% 16*f'(x)^2

% Sample points for f(x):
% x=-5: 25+3*sin^2(5)=25+3*(−0.9589)^2=25+2.758=27.758 -> y_t=4.72
% x=-4.5: 20.25+3*sin^2(4.5)=20.25+3*(−0.9775)^2=20.25+2.866=23.116 -> y_t=3.93
% x=-4: 16+3*sin^2(4)=16+3*(−0.7568)^2=16+1.718=17.718 -> y_t=3.01
% x=-3.5: 12.25+3*sin^2(3.5)=12.25+3*(−0.3508)^2=12.25+0.369=12.619 -> y_t=2.15
% x=-3: 9+3*sin^2(3)=9+3*(0.1411)^2=9+0.0597=9.060 -> y_t=1.54
% x=-2.5: 6.25+3*sin^2(2.5)=6.25+3*(0.5985)^2=6.25+1.074=7.324 -> y_t=1.245
% x=-2: 4+3*sin^2(2)=4+3*(0.9093)^2=4+2.480=6.480 -> y_t=1.102
% x=-1.5: 2.25+3*sin^2(1.5)=2.25+3*(0.9975)^2=2.25+2.985=5.235 -> y_t=0.890
% x=-1: 1+3*sin^2(1)=1+3*(0.8415)^2=1+2.124=3.124 -> y_t=0.531
% x=-0.5: 0.25+3*sin^2(0.5)=0.25+3*(0.4794)^2=0.25+0.690=0.940 -> y_t=0.160
% x=0:   0+0=0 -> y_t=0
% x=0.5: 0.940 -> y_t=0.160
% x=1:   3.124 -> y_t=0.531
% x=1.5: 5.235 -> y_t=0.890
% x=2:   6.480 -> y_t=1.102
% x=2.5: 7.324 -> y_t=1.245
% x=3:   9.060 -> y_t=1.54
% x=3.5: 12.619 -> y_t=2.15
% x=4:   17.718 -> y_t=3.01
% x=4.5: 23.116 -> y_t=3.93
% x=5:   27.758 -> y_t=4.72

\draw[stblue, line width=2.5pt] plot[smooth, tension=0.4]
  coordinates {
    (0,4.72) (0.7,3.93) (1.4,3.01) (2.1,2.15) (2.8,1.54)
    (3.5,1.245) (4.2,1.102) (4.9,0.890) (5.6,0.531) (6.3,0.160)
    (7.0,0) (7.7,0.160) (8.4,0.531) (9.1,0.890) (9.8,1.102)
    (10.5,1.245) (11.2,1.54) (11.9,2.15) (12.6,3.01) (13.3,3.93) (14,4.72)
  };

% 16 * f'(x)^2
% f'(x) = 2x + 3*sin(2x)
% x=-5: f'=−10+3*sin(−10)=−10+3*(0.5440)=−8.368, 16*70.02=1120 -> clipped
% x=-4: f'=−8+3*sin(−8)=−8+3*(−0.9894)=−10.968, 16*120.3=1924 -> clipped
% x=-3: f'=−6+3*sin(−6)=−6+3*(0.2794)=−5.162, 16*26.64=426 -> clipped
% x=-2.5: f'=−5+3*sin(−5)=−5+3*(0.9589)=−2.123, 16*4.51=72.1 -> clipped
% x=-2: f'=−4+3*sin(−4)=−4+3*(0.7568)=−1.729, 16*2.99=47.9 -> y_t=8.14
% x=-1.5: f'=−3+3*sin(−3)=−3+3*(−0.1411)=−3.423, 16*11.72=187.5 -> clipped
% x=-1: f'=−2+3*sin(−2)=−2+3*(−0.9093)=−4.728, 16*22.35=357.6 -> clipped

% The gradient bound grows fast so it will be mostly above the plot for large |x|.
% Let me focus on the region near 0 where f(x) and the bound are comparable.

% Near x=0: f'(0)=0, 16*0=0
% x=0.2: f'=0.4+3*sin(0.4)=0.4+3*(0.3894)=1.568, 16*2.459=39.3 -> y_t=6.68
% x=0.1: f'=0.2+3*sin(0.2)=0.2+3*(0.1987)=0.796, 16*0.634=10.1 -> y_t=1.72
% x=-0.1: same magnitude
% x=0.3: f'=0.6+3*sin(0.6)=0.6+3*(0.5646)=2.294, 16*5.26=84.2 -> clipped

% The bound 16*f'^2 is huge except very near 0. Let me clip at y=50 (tikz y=8.5)
% and plot within the visible range.

% Actually the Plotly code uses xrange [-5,5]. The gradient bound goes up fast.
% Let me sample densely near zero and clip appropriately.

% Manually computed points for 16*(2x+3*sin(2x))^2, clipped at 50:
% x=-0.6: f'=-1.2+3*sin(-1.2)=-1.2+3*(-0.9320)=-3.996, 16*15.97=255 -> clip
% x=-0.4: f'=-0.8+3*sin(-0.8)=-0.8+3*(-0.7174)=-2.952, 16*8.71=139.4 -> clip
% x=-0.25: f'=-0.5+3*sin(-0.5)=-0.5+3*(-0.4794)=-1.938, 16*3.76=60.1 -> clip
% x=-0.2: f'=-0.4+3*sin(-0.4)=-0.4-1.168=-1.568, 16*2.459=39.3 -> y_t=6.68
% x=-0.15: f'=-0.3+3*sin(-0.3)=-0.3-0.887=-1.187, 16*1.410=22.6 -> y_t=3.84
% x=-0.1: f'=-0.2+3*sin(-0.2)=-0.2-0.596=-0.796, 16*0.634=10.14 -> y_t=1.72
% x=-0.05: f'=-0.1-0.300=-0.400, 16*0.160=2.56 -> y_t=0.435
% x=0: 0 -> 0
% x=0.05: 2.56 -> y_t=0.435
% x=0.1: 10.14 -> y_t=1.72
% x=0.15: 22.6 -> y_t=3.84
% x=0.2: 39.3 -> y_t=6.68
% x=0.22: f'=0.44+3*sin(0.44)=0.44+1.278=1.718, 16*2.952=47.2 -> y_t=8.03

% The bound exits the plot around |x|=0.23.
% For the wider region, f(x) is always below the bound (PL condition).

\draw[terracotta, line width=2pt, dashed] plot[smooth, tension=0.4]
  coordinates {
    (6.08,8.5)
    (6.16,6.68)
    (6.37,3.84)
    (6.58,1.72)
    (6.79,0.435)
    (7.0,0)
    (7.21,0.435)
    (7.42,1.72)
    (7.63,3.84)
    (7.84,6.68)
    (7.92,8.5)
  };

% Add arrows to show the bound goes much higher outside
\draw[terracotta, line width=1.5pt, ->] (6.08,8.5) -- (5.8,8.8);
\draw[terracotta, line width=1.5pt, ->] (7.92,8.5) -- (8.2,8.8);

% Legend
\fill[cream, rounded corners=5pt] (8.5,6.0) rectangle (14.8,8.8);
\draw[border, line width=0.6pt, rounded corners=5pt] (8.5,6.0) rectangle (14.8,8.8);
\draw[stblue, line width=2.5pt] (8.9,8.2) -- (10.0,8.2);
\node[font=\large, anchor=west] at (10.2,8.2) {$f(x) = x^2\!+\!3\sin^2\!x$};
\draw[terracotta, line width=2pt, dashed] (8.9,7.2) -- (10.0,7.2);
\node[font=\large, anchor=west] at (10.2,7.2) {$16 \cdot |f'(x)|^2$};
\node[font=\large, text=wmgray, anchor=west] at (8.9,6.4) {PL: $f(x) \leq \tfrac{1}{2\mu}\|\nabla f\|^2$};

% Minimum marker
\fill[stblue] (7.0,0) circle (3.5pt);
\node[font=\large, text=stblue, below] at (7.0,-0.4) {$x^* = 0$};

% Annotation about nonconvexity
\draw[wmgray, rounded corners=5pt, line width=0.6pt] (2.0,-2.3) rectangle (12.5,-1.5);
\fill[wmgray, opacity=0.04, rounded corners=5pt] (2.0,-2.3) rectangle (12.5,-1.5);
\node[font=\large\itshape, text=dkbrown] at (7.25,-1.9) {Nonconvex, yet satisfies gradient domination with $\mu = 1/32$.};

\end{tikzpicture}
\end{document}
